\documentclass[11pt]{article}
\usepackage[margin=1in]{geometry}
\usepackage{amsmath,amssymb,amsthm}
\usepackage{booktabs}

\newtheorem{theorem}{Theorem}
\newtheorem{lemma}{Lemma}
\newtheorem{corollary}{Corollary}
\newcommand{\F}{\mathbb{F}}
\newcommand{\bd}{\partial}

\title{V72: Rank-Three Linear-Width Complexity and Branch Residual Composition}
\author{NC0\_k-Avoid Laboratory}
\date{31 July 2026}

\begin{document}
\maketitle

\begin{abstract}
We formalize three results for the support-hypergraph interface developed in
Laboratories V70--V71.  Vertex-boundary linear branch-width is NP-complete even
on simple three-uniform hypergraphs; an exact subset dynamic program runs in
$O(m2^m\operatorname{poly}(n))$ time; and a supplied binary gate decomposition
of boundary width $b$ supports exact affine residual composition using at most
$A(b)$ states per node and $A(b)^2$ child pairs per join.  We also record a
bounded-primal-treewidth family with unbounded linear support width.  None of
these statements resolves unrestricted Range Avoidance or P versus NP.
\end{abstract}

\section{Definitions}
Let $H=(V,E)$ be a hypergraph with labelled hyperedges.  For $S\subseteq E$,
define
\[
 \lambda_H(S)=\bigl|V(S)\cap V(E\setminus S)\bigr|.
\]
For an ordering $\pi=(e_1,\ldots,e_m)$, put
\[
 q(\pi)=\max_i\lambda_H(\{e_1,\ldots,e_i\}),\qquad
 q^*(H)=\min_\pi q(\pi).
\]
This is the V70 support-frontier optimum, expressed in the standard linear
branch-width vocabulary fixed in V71.

\section{Complexity at rank three}
\begin{theorem}[Three-uniform NP-completeness]
Given a simple three-uniform hypergraph $H$ and an integer $k$, deciding whether
$q^*(H)\le k$ is NP-complete.
\end{theorem}

\begin{proof}
An edge order is a polynomial-size certificate, and its maximum boundary can be
checked by maintaining processed and unprocessed incidence counts.

For hardness, start with a simple graph $G=(V,E)$.  For every edge $e=uv$,
introduce a fresh vertex $z_e$ and replace $e$ by the triple
$\{u,v,z_e\}$.  Call the resulting simple three-uniform hypergraph $H_G$.
For every $S\subseteq E$, no $z_e$ lies in the boundary because it occurs in
only one hyperedge.  An original vertex $v$ lies in the hypergraph boundary
exactly when it is incident to graph edges on both sides of the same cut.
Therefore
\[
 \lambda_{H_G}(S)=\lambda_G(S)
\]
for every labelled edge subset, so every ordering has exactly the same width in
$G$ and $H_G$.  Hence $q^*(H_G)$ equals graph linearwidth.  Computing graph
linearwidth is NP-hard, and the reduction is polynomial.
\end{proof}

\section{Exact subset optimization}
\begin{theorem}[Exact width dynamic program]
The value $q^*(H)$ and an optimal labelled-edge order can be computed in
$O(m2^m\operatorname{poly}(n))$ time using $O(2^m)$ numeric states.
\end{theorem}

\begin{proof}
For $S\subseteq E$, let $D[S]$ be the minimum maximum boundary among orders of
exactly the edges in $S$.  Then
\[
 D[\varnothing]=0,\qquad
 D[S]=\min_{e\in S}\max\{D[S\setminus\{e\}],\lambda_H(S)\}.
\]
In an optimal order of $S$, condition on its last edge for the lower bound; for
the upper bound, append that edge to an order realizing the predecessor state.
There are $2^m$ states and at most $m$ transitions per state.  Predecessor
pointers recover an optimal order.
\end{proof}

\section{Branch residual composition}
Consider an affine-cell gate system and a rooted binary tree whose leaves are
the gates.  For a node $t$, let $E_t$ be its leaf set and
\[
 \bd t=V(E_t)\cap V(E\setminus E_t).
\]
Let $R_t$ contain the distinct nonempty affine residuals obtained by selecting
one cell from each gate in $E_t$ and existentially projecting to $\bd t$.
Let $A(b)$ denote the number of nonempty affine subspaces of $\F_2^b$.

\begin{theorem}[Exact branch residual DP]
If $t$ has children $u,v$, then
\[
 R_t=\left\{
 \operatorname{proj}_{\bd t}(A\cap B):
 A\in R_u,\ B\in R_v,\ A\cap B\ne\varnothing
 \right\}.
\]
Consequently, on a supplied decomposition with $|\bd t|\le b$ for every node,
exact affine feasibility is computable in
\[
 O\!\left(mA(b)^2\operatorname{poly}(n)\right)
\]
time with at most $A(b)$ residual states per node.
\end{theorem}

\begin{proof}
Induct on the binary tree.  At a leaf, project its two cells and deduplicate.
At an internal node, every cell selection splits uniquely into selections in
the two child subtrees.  A variable shared by the children is visible outside
each child, so child projection cannot discard information required by their
conjunction.  The conjunction of the two child residuals followed by projection
to $\bd t$ is therefore exactly the direct residual of the combined selection.
The reverse inclusion follows from the child selections witnessing each pair.

A residual on at most $b$ Boolean coordinates is a nonempty affine subspace of
$\F_2^b$, giving $|R_t|\le A(b)$.  Each join tests at most $A(b)^2$ pairs, and
a binary tree with $m$ leaves has $O(m)$ nodes.
\end{proof}

\section{Treewidth does not control the linear parameter}
\begin{corollary}
There is a family of simple three-uniform hypergraphs whose primal graphs have
treewidth at most two while $q^*$ is unbounded.
\end{corollary}

\begin{proof}
Apply the private-vertex construction to perfect binary trees.  A tree edge and
its private vertex form a triangle, and one three-vertex bag per original edge
gives a width-two tree decomposition.  The construction preserves every
edge-cut boundary, so its $q^*$ equals the graph linearwidth of the tree.
Perfect binary trees have unbounded pathwidth, and graph linearwidth differs
from pathwidth by at most one outside the known tiny exception.  Thus $q^*$ is
unbounded.
\end{proof}

\section{Finite benchmark}
On six frozen V69--V70 systems, pathwidth-derived orders have projected-DAG
cost between $1.379$ and $2.059$ times the exact preserved $G^*_{\rm proj}$.
This is finite experimental evidence, not an approximation theorem.  The exact
width DP is independently checked against permutation enumeration, and branch
residual composition is checked against direct cell enumeration.

\section{Scope}
The branch computation is tree-shaped affine feasibility, not the linear
$G_{\rm proj}$ model and not an asserted OBDD, FBDD, resolution, Res-Lin, or
communication-complexity simulation.  We do not prove a general polynomial
good-order theorem, an all-orders superpolynomial lower bound, unrestricted
$\mathrm{NC}^0_3$-Avoid, or any P-versus-NP consequence.  Novelty and peer
review are unconfirmed.

\end{document}
